\scanpage{205}
% archive sheet number: 204
\deleted{Il faut \unsure{donc} \unsure{prendre} $E_3^{pq}$}

Il faut \unsure{croire|voir} que $E_2^{p*}\longrightarrow$ \inserted{(\unsure{un. autre twist|un autre twist|ou autres twists})} le \deleted{\unsure{complexe des}} complexe dual de \deleted{$\unsure{E_2^{**}(Y)|E_2^{**}}$} $\sigma\mapsto H^p(Y_\sigma)$ \inserted{de \textcircled{a}}, \unsure{compte} \unsure{tenu} de la dualité de Poincaré, de sorte que \deleted{\unsure{var. de|via} $\unsure{Y_\sigma}$} ${}^{b}E_3^{**}$ est dual de ${}^{(a)}E_2^{**}$ à twist près.

\textcircled{c} Autre présentation \inserted{de $H^{\unsure{*|\bullet}}(X,\mathbb{Q}_{\unsure{\ell|l}})$} : suite spectrale de Leray de $\unsure{g|j}:X\hookrightarrow X'$; \unsure{=}
\[
H^*(X)\Longleftarrow E_2^{pq}=H^p(X',R^q g_*(\mathbb{Q}_\ell))
\]
Or
\[
\begin{aligned}
R^q g_*(\mathbb{Q}_\ell)&\simeq\mathop{\bigwedge}^{q}\coprod_\alpha\mathbb{Q}_\ell(-1)_{Y_\alpha}\\
&\simeq\coprod_{\alpha_1<\cdots<\alpha_q}\mathbb{Q}_{\ell Y_{\alpha_1\cdots\alpha_q}}(-q).
\end{aligned}
\]
Donc
\[
\left\{\begin{aligned}
E_2^{pq}&=\coprod_{\alpha_1<\cdots<\alpha_q}H^p(Y_{\alpha_1\cdots\alpha_q})(-q)\qquad\text{si }q\geq1\\
E_2^{p0}&=H^p(X')\deleted{\unsure{\mathbb{Q}_\ell}}
\end{aligned}\right.
\]
qui est \uline{de poids} $p+2q$. \uline{Donc on trouve} \deleted{\illegible{3}}
\[
\underline{d_r=0\quad\text{si }r\geq3},
\]
et il faut calculer encore $E_3^{pq}$, \deleted{\unsure{comme}} i.e. les
\diagram{\begin{tikzcd}[ampersand replacement=\&,column sep=small,row sep=large]
\&E_2^{pq}\arrow[r]\arrow[dl,equal]\&E_2^{p+2,q-1}\arrow[dr,equal]\&\\
\displaystyle\coprod_{\alpha_1<\cdots<\alpha_q}H^p(\tikz[baseline=(yl.base)]{\node[inner sep=0pt] (yl) {$Y_{\alpha_1\cdots\alpha_q}$};\node[anchor=east] (dl) at (-.55,.8) {\uline{dim $\unsure{n-1-q}$}};\draw (dl.east) to[out=0,in=90] (yl.north);})(-q)\arrow[rrr]\arrow[d,dash,"{\text{dual}}"']\&\&\&\displaystyle\coprod_{\alpha_1<\cdots<\alpha_{q-1}}H^{p+2}(\tikz[baseline=(yr.base)]{\node[inner sep=0pt] (yr) {$Y_{\alpha_1\cdots\alpha_{\unsure{q-3|q-1}}}$};\node[anchor=west] (dr) at (.55,.8) {dim $n-q$};\draw (dr.west) to[out=180,in=90] (yr.north);})(-(q-1))\arrow[d,dash,"{\text{dual}}"]\\
H^{\unsure{2n-2-p-2q|2n-2-p+2q}}(Y_{\alpha_1\cdots\alpha_q})(n-1)\&\&\&H^{2n-2q-p-2}(Y_{\alpha_1\cdots\alpha_{q-1}})(n-1)
\end{tikzcd}}
% Dimension notes have pencil leaders to the respective Y terms; the lower "dual" lines have no arrowheads.
\margin{si $q\geq2$\\i.e. $q-1\geq1$}
\figure{\begin{tikzpicture}\draw (0,0)--(0,3);\end{tikzpicture}}
